\documentclass{article}
\usepackage{graphicx} % Required for inserting images
\usepackage{caption} % Required for \captionof
\usepackage{amsmath}
\usepackage{float}
\usepackage{comment}
\usepackage[T1]{fontenc}
\usepackage[utf8]{inputenc}
\usepackage{lmodern}
\usepackage{microtype}
\usepackage{xcolor}
\usepackage{framed}
\renewenvironment{leftbar}{%
  \def\FrameCommand{\textcolor{orange}{\vrule width 3pt} \hspace{10pt}}%
  \MakeFramed {\advance\hsize-\width \FrameRestore}}%
 {\endMakeFramed}

\title{Mechanikai mérés}
\author{Kern Luca, Vincze Csongor}
\date{2026 Április 22}

\begin{document}

\maketitle

\section*{2. Feladat: Határozza meg a rendszer $\beta = \frac{\gamma}{2}$ csillapítási tényezőjét!}

\begin{table}[H]
    \centering
    \begin{tabular}{|c|c|c|}
        \hline
        Szögkitérítés [°] & Idő [s] & Periódusidő [s] \\ \hline
        90 & 12,81 & \\ \hline 
        100 & 13,06 & \\ \hline
        110 & 12,85 & \\ \hline 
        120 & 12,84 & \\ \hline
        130 & 12,83 & \\ \hline 
        140 & 12,81 & \\ \hline
        150 & 12,79& \\ \hline 
         &  & \\ \hline
        
    \end{tabular}
    \caption{Periódusidő a kezdeti szögkitérítéstől függően}
    \label{tab:periodusido}
\end{table}

20 lengés után nézzük a maximális amplitúdót.
$\phi_0 = 90^\circ$
%videó!!!!!

\begin{table}[H]
    \centering
    \begin{tabular}{|c|c|}
        \hline
        Mérések száma [db] & Maximális amplitúdó [°] \\ \hline
         & \\ \hline 
         & \\ \hline
         & \\ \hline 
         & \\ \hline
         & \\ \hline 
         & \\ \hline
         & \\ \hline 
         & \\ \hline
        
    \end{tabular}
    \caption{Maximális amplitúdó}
    \label{tab:max_amplitudo}
\end{table}



\end{document}
